\clearpage
\begin{trapbox}
\textbf{误区一：比例取倒数，或把平分角当成平分边。}
错误写成 $BD/DC=AC/AB$，或不论邻边就写 $BD=DC$。
下图固定 $AB=6,AC=4,BC=5$，正确分点为 $D$，倒置比例产生错误点 $D_*$。
\begin{center}
\begin{tikzpicture}[font=\small]
\draw[black!30,rounded corners=2pt] (0,0) rectangle (14.4,8.2);
\node[anchor=west,font=\small\bfseries] at (0.5,7.72) {倒置比例会换分点：同一三角形中的 $D$ 与错误点 $D_*$};
% MODEL 6 4 5
\begin{scope}[shift={(4,2)},scale=1.15]
\coordinate (A) at (4.5,3.9686269666);
\coordinate (B) at (0,0);
\coordinate (C) at (5,0);
\coordinate (D) at (3,0);
\coordinate (W) at (2,0);
\draw[blue,thick] (A)--(B);\draw[teal!80!black,thick] (A)--(C);\draw[thick] (B)--(C);\draw[orange!85!black,very thick] (A)--(D);
\fill (A) circle (1.3pt) node[above=5pt] {$A$};
\fill (B) circle (1.3pt) node[below left=5pt] {$B$};
\fill (C) circle (1.3pt) node[below right=5pt] {$C$};
\fill (D) circle (1.3pt) node[below=5pt] {$D$};
\draw[red!75!black,thick,dashed] (A)--(W);\fill[red!75!black] (W) circle (1.3pt) node[below left=6pt] {$D_*$};
\path (A)--(B) node[midway,above left=4pt] {$AB=6$};
\path (A)--(C) node[midway,above right=4pt] {$AC=4$};
\end{scope}
\node[anchor=west,text=orange!85!black] at (.8,6.5) {正确点：$3:2$};
\node[anchor=west,text=red!75!black] at (.8,5.7) {倒置后：$2:3$};
\draw[blue,|-|] ($(B)+(0,-1.05)$)--($(D)+(0,-1.05)$) node[midway,below=3pt] {$BD=3$};
\draw[teal!80!black,|-|] ($(D)+(0,-1.05)$)--($(C)+(0,-1.05)$) node[midway,below=3pt] {$DC=2$};

\end{tikzpicture}
\end{center}

正确替代：按 $B\leftrightarrow AB$、$C\leftrightarrow AC$ 对应。
本图 $BD=3,DC=2$；错点满足 $BD_*/D_*C=2/3$，不是正确的 $3/2$。
若误设等长，则取到中点；它与真正角平分线的差别已由“角平分线、中线、高”图直接展示。

\tcbbreak
\textbf{误区二：逆用只看比例，不看点的位置。}
错误做法：只要 $BX/XC=AB/AC$，就断言 $AX$ 平分内角。
\begin{center}
\begin{tikzpicture}[font=\small]
\draw[black!30,rounded corners=2pt] (0,0) rectangle (14.4,8.2);
\node[anchor=west,font=\small\bfseries] at (0.5,7.72) {逆定理不能漏掉位置：内点 $D$ 与外点 $X$ 的长度比相同};
% MODEL 6 4 5
\begin{scope}[shift={(2.1,2.5)},scale=0.68]
\coordinate (A) at (4.5,3.9686269666);
\coordinate (B) at (0,0);
\coordinate (C) at (5,0);
\coordinate (D) at (3,0);
\coordinate (X) at (15,0);
\draw[blue,thick] (A)--(B);\draw[teal!80!black,thick] (A)--(C);\draw[thick] (B)--(C);\draw[orange!85!black,very thick] (A)--(D);
\fill (A) circle (1.3pt) node[above=5pt] {$A$};
\fill (B) circle (1.3pt) node[below left=5pt] {$B$};
\fill (C) circle (1.3pt) node[below right=5pt] {$C$};
\fill (D) circle (1.3pt) node[below=5pt] {$D$};
\draw[black!45] (C)--(X);\draw[purple,thick,dashed] (A)--(X);\fill[purple] (X) circle (1.3pt) node[below right=5pt] {$X$};
\path (A)--(B) node[midway,above left=4pt] {$AB=6$};
\path (A)--(C) node[midway,above right=4pt] {$AC=4$};
\path (C)--(X) node[midway,above=6pt] {$XC=10$};
\end{scope}
\draw[blue,|-|] ($(B)+(0,-1.1)$)--($(D)+(0,-1.1)$) node[midway,below=3pt] {$BD=3$};
\draw[teal!80!black,|-|] ($(D)+(0,-1.1)$)--($(C)+(0,-1.1)$) node[midway,below=3pt] {$DC=2$};
\node[anchor=west,text=purple] at (8.1,6.4) {$BX/XC=15/10=3/2$};
\node[anchor=west,text=purple] at (8.1,5.6) {$X$ 在外部，$AX$ 平分外角};

\end{tikzpicture}
\end{center}

图中内点 $D$ 满足 $BD/DC=3/2$；外点 $X$ 也满足 $BX/XC=15/10=3/2$。
但是 $AX$ 是外角平分线，不能当作内角平分线。
正确替代：逆定理先检查 $D$ 严格位于 $BC$ 内部，再检查比例。
内部用 $BD+DC=BC$，外部用长度差；本条只用前一种。

\textbf{误区三：条件不成立，仍机械解方程。}
原例 $9,6,15$ 不满足严格三角形不等式，不能凭算出分段比就说三角形存在；例 1 的退化对照图已画出三点共线。
正确替代：先查 $|AB-AC|<BC<AB+AC$，再计算。

\textbf{误区四：逆证明中直接把正弦相等当成角相等。}
$0^\circ$ 到 $180^\circ$ 内，正弦相等还可能来自补角。
若走三角函数路线，须再用 $\angle BAD+\angle DAC<180^\circ$ 排除补角；
本条的逆证明使用唯一分点，避开了这个分支遗漏。
\end{trapbox}
